\documentclass{crypto-exercise}
\author{Sven Laur}
\editor{Sven Laur}

\tags{simple reductions, Discrete Logarithm problem, incomplete
  solution, star exercise}

\begin{document}
\begin{exercise}{Hardness of intermediate discrete logarithm bits}
  Let $\GG=\langle g\rangle$ be a finite group of an odd order $q$
  generated by the powers of an element $g$. Then the Discrete
  Logarithm (DL) problem is following. For any element $y$ find a
  power $x\in\ZZ_q$ such that $g^x=y$. Let $\AD$ be an algorithm that computes the $\ell$-th bit of discrete logarithm, i.e., the algorithm is correct if $\AD(g^{x_n\ldots x_1})=x_\ell$. 
  
  
  \begin{enumerate}
  \item Show that if $\AD$ is always correct then there exists an efficient procedure that recovers the discrete logarithm with probability $2^{\ell-1}$.
  \item Show that if $\AD$ is always correct then there exists an efficient procedure that always recovers the discrete logarithm but runs $2^{\ell-1}$ times slower.
  \item   Let $\hat{y}_1=\AD(y),\hat{y}_2=\AD(y^2),\ldots, \hat{y}_n=\AD(y^{2^{n-1}})$ for an algorithm $\AD$ that is always correct. 
  Write down the corresponding system of linear equations over the integers for the unknowns 
$x_1,\ldots,x_n$, together with the modular-reduction correction terms. 
These correction terms arise from the fact that exponents are always taken modulo $q$. Is it possible to solve this equation?
\item[($\star$)] The integer equation formalised above has exponentially many solutions and only for one of them $x_n,\ldots,x_1\in\set{0,1}$. 
Note that the nullspace over integers is a lattice. Formalise recovery of  $x_n,\ldots,x_1$ as a closest vector problem. Does this reduces the runningtime estmate for the algorithm that always recovers $x$.
 \end{enumerate} 
\end{exercise}


\begin{solution}
\end{solution}

\end{document}
